\chapter{Introduction}
\label{cap:introduction}

\section{Overview}

Technology has always influenced society, from the first introduction to the Personal Computer to the Internet, changing the world year after year, from personal life to workspaces, people are surrounded by software on their computers and mobile devices. Over the last decades, the investments in the development of technological innovations, both in hardware and software, as increased exponentially, shaping economical and social impact by integrating the latest features into everyday life. Smartphones and tablets have been the core engine that completely changed daily habits and digital usage; their wide adoption has reached almost three quarters of global population \cite{site:mobile-ownership,site:mobile-sales}, making users abandon traditional computers, not just for web browsing but more especially for social media\cite{site:mobile-facebook-usage} and instant messaging platform\cite{site:mobile-messageapp-downloads}, that changed the design and development of applications into a `mobile first' logic.


Compared to computers, mobile devices offer distinct functionalities and strengths that have facilitated their widespread adoption. These can be categorized according to the fundamental logic underlying any technology:
\begin{enumerate}
    \item \textbf{Hardware}: the increase in computational performance\cite{site:mobile-performance}, battery life\cite{site:mobile-battery-increase}, fast connectivity, and the integration of several built-in sensors\cite{site:mobile-sensors} such as gyroscopes, GPS, NFC, light sensors, cameras, and vibration motors. These enable applications to interact with different fields such as health\cite{site:mobile-health-monitoring}, education\cite{site:mobile-education}, and, through motion-sensors or cameras, an interaction with physical world, creating more immersive user experience.

    \item \textbf{Software}: the development and presence of mobile applications has played a crucial role. In fact, platforms such as `Windows Phone' were eventually abandoned and discontinued due to the lack of applications \cite{site:windowsphone-appgap-1,site:windowsphone-appgap-2} compared to operating systems like `iOS'\cite{site:iOS} and `Android'\cite{site:Android}. 
\end{enumerate}


In this way, smartphones apps have became a key integration with daily life eand physical world, for example, when planning a trip or just driving, smartphones, integrated with cars assists users via app like `Google Maps'\cite{site:google-maps} or `Waze'\cite{site:waze}, providing not just directions but also real-time traffic updates, read events or information about available parking. This makes the surrounding environment enriched with digital details. Another example is `Google Lens'\cite{site:google-lens}, which can identify objects from a photograph, or `Shazam'\cite{site:shazam} that through devices` microphone can recognize a playing song.
Even in workplaces or public areas, technology integrations can be present, from the resturant that provide a digital menu by a tablet, to customers checking the nutrition macros into a grocery store by scanning the barcode on the pack, to museums by scanning a QR Code enabling audio-guided tours or unlock Augmented Reality functionalities.



\section{Social networks}




% Applications such as WhatsApp, Telegram, Instagram, and Facebook have attracted billions of users worldwide, becoming deeply embedded in daily communication practices across diverse social, professional, and public contexts. While gaining huge amount of users, these apps has also raised concerns about social influence, data privacy (which often isnt respected) and evereal other problems as shadow banning, fake news, cyber-bullism and others illegal content.
% These platforms have totally changed the way people interact among themself digitally, by providing instant-time messaging capabilities, both one-to-one as one-to-many communication and also enabling users to share any kind of multimedia (including photos, videos, audio recordings and other files). The content shared, usually referred as `post` crate a starting point for interactions between users, throught comments and reactions (i.e. `like',`dislike'), thus to digitally replicate human behavior during a dialogs and interactions.
% Despite positive effects of ocial media like connections and information sharing, some serious concerns have raised like data privacy (which often isn't properly respected from the platform in first place), but also shadow banning, fake news, cyber-bullisms, stalking and others illegal content. Platforms in order to safeguard users, has applied different mechanisms thust to make socials safer as possible, even if not always effective. Human reports / heading behavior remains the only true system to recognize not allowed content


Alongside the widespread adoption and continuous advancement of smartphones technologies and other mobile devices, social networks, have emerged as one of the most influential and important sectors of the technology market. Today, social media platforms rank among the firsts widely used applications worldwide \cite{site:social-media-usage-adults,site:social-media-usage-usa}, while the companies that operate them become some of the most bigger actors in global financial markets \cite{site:social-media-financial-market,site:social-media-financial-influence,site:social-media-meta-revenue}.
\newline


Social media can be divided into different subgroup, like instant-messaging or online-social-network, but despite their nature they all provide some key points like instant-time messaging capabilities, both one-to-ton as one-to-many communications till to be integrated in 


Platform like WhatsApp, Telegram, Instagram, and Facebook have totally changed the way people interact among themself digitally, by providing instant-time messaging capabilities, both one-to-one as one-to-many communication and also enabling users to share any kind of multimedia (including photos, videos, audio recordings and other files) till to be integrated into everyday communication practices across social, professional, and public environments.

The content shared, usually referred as `post` crate a starting point for interactions between users, throught comments and reactions (i.e. `like',`dislike'), thus to digitally replicate human behavior during a dialogs and interactions





In addition to messaging capabilities, these platforms allow users to create and share a wide range of multimedia content, including images, videos, audio recordings, and documents. Such content, commonly referred to as a *post*, serves as the basis for user engagement through comments, reactions, and content redistribution mechanisms. These features digitally replicate several aspects of human social interaction, fostering the creation of large online communities and enabling information to spread rapidly across extensive networks of users.

Despite the significant benefits offered by social media platforms, including enhanced connectivity, community building, and information sharing, their widespread adoption has also introduced a number of societal and ethical challenges. Concerns regarding user privacy, the collection and monetization of personal data, the dissemination of misinformation and fake news, cyberbullying, stalking, online harassment, and the circulation of illegal or harmful content have become increasingly relevant topics of discussion among researchers, policymakers, and the general public.

To address these issues, social media platforms have implemented various moderation and safety mechanisms, including automated content analysis systems, community guidelines, reporting tools, and human review processes. While these measures contribute to reducing the spread of harmful content, they are not always fully effective due to the scale and complexity of modern social networks. Consequently, user reporting and human moderation continue to play a crucial role in identifying, evaluating, and removing content that violates platform policies or legal regulations.






% ------------------------------------------------------


% \noindent Esempio di utilizzo di un termine nel glossario \\
% \gls{api}. \\

% \noindent Esempio di citazione in linea \\
% \cite{site:agile-manifesto}. \\

% \noindent Esempio di citazione nel pie' di pagina \\
% citazione\footcite{womak:lean-thinking} \\

% \section{L'azienda}

% Descrizione dell'azienda.

% \section{L'idea}

% Introduzione all'idea dello stage.

% \section{Organizzazione del testo}

% \begin{description}
%     \item[{\hyperref[cap:processi-metodologie]{Il secondo capitolo}}] descrive ...
    
%     \item[{\hyperref[cap:descrizione-stage]{Il terzo capitolo}}] approfondisce ...
    
%     \item[{\hyperref[cap:analisi-requisiti]{Il quarto capitolo}}] approfondisce ...
    
%     \item[{\hyperref[cap:progettazione-codifica]{Il quinto capitolo}}] approfondisce ...
    
%     \item[{\hyperref[cap:verifica-validazione]{Il sesto capitolo}}] approfondisce ...
    
%     \item[{\hyperref[cap:conclusioni]{Nel settimo capitolo}}] descrive ...

%     \begin{lstlisting}[language=Python, caption={Python example}]
%     def hello():
%         print("Hello")
%     \end{lstlisting}

%     \begin{lstlisting}[language=JavaScript, caption={JavaScript example}]
%     function hello() {
%         console.log("Hello");
%     }
%     \end{lstlisting}

%     \begin{lstlisting}[language=SQL, caption={SQL example}]
%     SELECT * FROM users;
%     \end{lstlisting}

% \end{description}



% Riguardo la stesura del testo, relativamente al documento sono state adottate le seguenti convenzioni tipografiche:
% \begin{itemize}
% 	\item gli acronimi, le abbreviazioni e i termini ambigui o di uso non comune menzionati vengono definiti nel glossario, situato alla fine del presente documento;
% 	\item per la prima occorrenza dei termini riportati nel glossario viene utilizzata la seguente nomenclatura: \emph{parola}\glsfirstoccur;
% 	\item i termini in lingua straniera o facenti parti del gergo tecnico sono evidenziati con il carattere \emph{corsivo}.
% \end{itemize}
